Deep reinforcement learning (DRL) has emerged as a transformative tool in the modeling, control, and maintenance of power electronic converter systems. This paper provides a comprehensive review of the current state-of-the-art in applying DRL techniques across the converter system lifecycle, covering design-stage modeling, control strategies, and condition-based maintenance.

We first reviewed recent progress in using DRL for intelligent control of power converters, categorizing applications into five representative domains: maximum power point tracking (MPPT), modulation strategy optimization, variable-parameter tuning, energy management, and motor speed and position tracking. In each domain, DRL demonstrates strong capabilities in adapting to nonlinearities, uncertainties, and time-varying dynamics, outperforming many traditional rule-based and model-dependent control methods.

In the context of converter system maintenance, we further reviewed how DRL is being used for condition monitoring, anomaly detection and fault diagnosis, and remaining useful life (RUL) prediction. The reviewed works highlight DRL’s potential to realize scalable, autonomous, and data-driven health management systems through techniques such as attention-based sequence modeling, actor-critic frameworks, and hybrid reinforcement learning architectures.

Despite these advancements, our review also reveals critical limitations that hinder practical deployment, including simulation-heavy validation, lack of policy interpretability, computational constraints, and the absence of standardized DRL deployment frameworks in embedded power hardware. To address these gaps, we outlined future research directions such as hardware-aware DRL implementations, physics-informed learning, explainable DRL, and multi-agent coordination.

In conclusion, DRL holds considerable promise in enabling the next generation of intelligent, autonomous, and resilient power electronic converter systems. However, realizing this vision will require continued interdisciplinary efforts in algorithm development, system integration, and real-time implementation. Bridging the gap between simulation and deployment, while ensuring safety, interpretability, and efficiency, will be essential for DRL to reach its full potential in the power electronics domain.
